\documentclass[../main.tex]{subfiles}\begin{document}

% UI/UX

\section{Frontend}
Dette afsnit analyserer valg af frontend-teknologi for Echo. Gruppen har valgt at udvikle løsningen som en web-app frem for en separat mobilapplikation, da en web-app kan bruges på tværs af desktop, tablet og mobil, hvis brugergrænsefladen designes responsivt. Dette passer godt til Echo, hvor Contributors, Requesters, Organizations og Reward Partners skal kunne tilgå platformen fra forskellige enheder uden at gruppen skal vedligeholde flere separate klienter.\\

Analysen fokuserer derfor på tre relevante frontend-frameworks til web-apps: Angular, React og Vue. Teknologierne vurderes ud fra både generelle kriterier og krav, der er særligt relevante for Echo, herunder formularer, rollebaseret navigation, API-integration, responsivt design og testbarhed.

\subsection{Krav}
Frontend-løsningen skal understøtte Echos centrale brugerflows på en intuitiv måde. Brugere skal kunne oprette konto, logge ind, vedligeholde profil, gennemse opgaver, acceptere eller tildele opgaver, se saldo og transaktionshistorik, indløse rewards og afgive ratings efter fuldførte opgaver. Da systemet har flere brugerroller, skal frontend også kunne håndtere forskellig navigation og visning afhængigt af om brugeren er Contributor, Requester, Organization eller Reward Partner.\\

De ikke-funktionelle krav stiller også krav til frontend-valget. Siden bør indlæses inden for 5 sekunder under normale netværksforhold, og løsningen bør understøtte mobilvisning. Derudover skal brugergrænsefladen være intuitiv for nye brugere, da platformens værdi afhænger af, at brugerne hurtigt kan forstå, hvordan de finder opgaver, håndterer belønninger og vurderer andre brugeres omdømme.\\

% Kriterier
For at analysere frontend-teknologierne er der opstillet kriterier, som både dækker generelle teknologivalg og de krav, der er specifikke for Echo.

% Kræver biblatex (eller natbib) og kilder.bib

% Kræver biblatex (eller natbib) og kilder.bib

\begin{table}[H]
    \centering
    \renewcommand{\arraystretch}{1.25}
    \begin{tabularx}{\textwidth}{|X|c|X|}
        \hline
        \rowcolor{gray!30}
        \textbf{Kriterie} & \textbf{Vægt} & \textbf{Relevans for Echo} \\
        \hline
        Egnethed til projektet & 20\% & Frameworket skal passe til en rollebaseret platform med opgaver, rewards, betaling og ratings. \\
        \hline
        Maintainability & 10\% & Frontenden skal kunne udvides uden at brugerflows og komponenter bliver svære at overskue. \\
        \hline
        Performance & 5\% & Web-appen skal kunne indlæses hurtigt og føles responsiv på både desktop og mobil. \\
        \hline
        Kompetencer i gruppen & 10\% & Valget skal afspejle gruppens erfaring med Angular, React og Vue. \\
        \hline
        Ecosystem/tooling & 5\% & Frameworket bør have gode værktøjer til routing, komponenter, test, linting og build. \\
        \hline
        Fremtidig skalerbarhed & 5\% & Frontenden skal kunne vokse med flere sider, roller og funktioner. \\
        \hline
        Responsivt webdesign og mobilvisning & 10\% & Echo skal kunne bruges på forskellige skærmstørrelser uden separate apps. \\
        \hline
        Formularer, routing og rollebaseret navigation & 10\% & Konto, login, oprettelse af opgaver og rewards kræver stærke formular- og routingmuligheder. \\
        \hline
        State management og API-integration & 10\% & Frontenden skal håndtere loginstatus, brugerroller, opgavelister, saldo og transaktionsdata. \\
        \hline
        Type safety og testbarhed & 10\% & Koden skal være robust og kunne testes, især ved flows med betaling, roller og opgavestatus. \\
        \hline
        Komponentstruktur og design system & 5\% & Genbrugelige komponenter er vigtige for ensartet UI på tværs af opgaver, rewards, profiler og ratings. \\
        \hline
        \rowcolor{gray!30}
        \textbf{Total} & \textbf{100\%} & \\
        \hline
    \end{tabularx}
    \caption{Vurderingskriterier for valg af frontend-teknologi}
    \label{tab:frontend-kriterier}
\end{table}

\subsection*{Egnethed til projektet (20\%)}
De tre kandidater adskiller sig mest i, hvor meget de leverer ud af boksen. Angular er et komplet framework med indbygget routing, formularhåndtering, HTTP-klient og dependency injection~\cite{angular-di}. Det passer til Echo, hvor mange rollebaserede flows skal struktureres ensartet. Vue er et progressivt framework med officielle løsninger til routing og state management i form af Vue Router og Pinia~\cite{vue-router-guards, pinia}, men færre indbyggede værktøjer end Angular, f.eks. til formularvalidering. React er et UI-bibliotek, hvor routing, formularer og datahentning kræver tredjepartsbiblioteker, og React-teamet anbefaler i dag at starte nye projekter med et framework oven på React~\cite{react-cra-sunset}. Både React og Vue kan bruges til Echo, men kræver, at gruppen selv træffer flere arkitekturvalg.

\medskip\noindent\textbf{Score:} Angular: 9 \quad React: 7 \quad Vue: 7

\subsection*{Maintainability (10\%)}
Angular er det mest meningsbærende af de tre. Opdelingen i komponenter, services og dependency injection giver en ensartet struktur, som gør koden lettere at overskue på tværs af gruppen. Angular udgiver en ny hovedversion hvert halve år med 18 måneders support, og \texttt{ng update} automatiserer mange migreringer~\cite{angular-releases}. Vue samler template, logik og styling i Single-File Components, hvilket gør den enkelte komponent overskuelig~\cite{vue-sfc}, men giver færre konventioner for strukturen af applikationen som helhed. React giver størst frihed, men netop derfor afhænger vedligeholdelsen af, at gruppen selv fastlægger og overholder konventioner for filstruktur og dataflow.

\medskip\noindent\textbf{Score:} Angular: 9 \quad React: 7 \quad Vue: 7

\subsection*{Performance (5\%)}
Alle tre frameworks er modne og optimerede, og js-framework-benchmark viser, at forskellene mellem dem er små i forhold til, hvad en bruger vil kunne mærke i en typisk forretningsapplikation~\cite{js-framework-benchmark}. Angular har tidligere været kritiseret for større bundles, men fra version 21 bruger nye projekter zoneless change detection, som ifølge Angular-teamet reducerer bundlestørrelsen og forbedrer Core Web Vitals~\cite{angular-v21}. Alle tre understøtter lazy loading af ruter. For Echo vil performance primært afhænge af API-kald og billeder frem for valget af framework.

\medskip\noindent\textbf{Score:} Angular: 8 \quad React: 8 \quad Vue: 8

\subsection*{Kompetencer i gruppen (10\%)}
Gruppen føler sig generelt ikke stærk i frontend-udvikling. Under praktikperioden og de efterfølgende studiejob har fokus primært ligget på backend, og i tidligere kurser har gruppen i vid udstrækning brugt AI-værktøjer til at udvikle frontenden, fordi andre dele af projekterne havde højere prioritet. Den praktiske erfaring med frontend-frameworks er derfor mindre, end antallet af gennemførte projekter antyder. React er det framework, gruppen har mest erfaring med, da det er blevet undervist i et kursus og anvendt i et tidligere semesterprojekt. Erfaringen med Angular er mere begrænset, og Angulars stejlere læringskurve gør oplæringen mere krævende. Vue har gruppen næsten ingen erfaring med. Uanset valg vil en del af projektets første fase skulle bruges på oplæring, hvilket er en risiko i et tidsbegrænset bachelorprojekt.

\medskip\noindent\textbf{Score:} Angular: 4 \quad React: 5 \quad Vue: 2

\subsection*{Ecosystem/tooling (5\%)}
React har det største ecosystem og bruges af knap halvdelen af respondenterne i Stack Overflows udviklerundersøgelse fra 2025, mens Angular og Vue hver bruges af under en femtedel~\cite{so-survey-2025}. Det giver React det største udvalg af biblioteker og det største community. Angular har til gengæld det mest integrerede tooling: Angular CLI håndterer generering af komponenter, build og test~\cite{angular-cli}, og siden version 21 er Vitest standard-testrunner~\cite{angular-v21}. Vue har et solidt, men mindre ecosystem, og bygger typisk på Vite og Vitest~\cite{vue-testing}.

\medskip\noindent\textbf{Score:} Angular: 9 \quad React: 10 \quad Vue: 8

\subsection*{Fremtidig skalerbarhed (5\%)}
Alle tre understøtter opdeling i komponenter og lazy loading, så applikationen kan vokse med flere sider og roller uden at blive tungere at indlæse. Angulars fastlagte struktur skalerer godt, når både kodebase og antal udviklere vokser, og frameworket bruges bredt i store virksomhedsapplikationer. React bruges ligeledes i meget stor skala, men forudsætter, at teamet har etableret faste konventioner. Vue skalerer også fint, men har færre udbredte konventioner for meget store applikationer.

\medskip\noindent\textbf{Score:} Angular: 9 \quad React: 9 \quad Vue: 8

\subsection*{Responsivt webdesign og mobilvisning (10\%)}
Responsivt design løses primært med CSS i form af media queries, flexbox og grid og er dermed i vid udstrækning uafhængigt af frameworket~\cite{mdn-responsive}. Alle tre kan kombineres med CSS-frameworks som Tailwind eller Bootstrap og har komponentbiblioteker med responsive komponenter, f.eks. Angular Material, MUI og Vuetify. Angular CDK tilbyder desuden \texttt{BreakpointObserver}, som kan tilpasse komponenters adfærd efter skærmstørrelse~\cite{angular-cdk-layout}. Kriteriet adskiller derfor ikke kandidaterne, men er medtaget, fordi mobilvisning er et centralt krav til Echo.

\medskip\noindent\textbf{Score:} Angular: 8 \quad React: 8 \quad Vue: 8

\subsection*{Formularer, routing og rollebaseret navigation (10\%)}
Echo indeholder mange formularer til konto, login, opgaver og rewards samt navigation, der afhænger af brugerens rolle. Angular har Reactive Forms med typede formularer og validering indbygget~\cite{angular-reactive-forms}, og routeren understøtter guards, der kan begrænse adgangen til ruter ud fra brugerens rolle~\cite{angular-route-guards}. Vue Router er officielt og understøtter navigation guards~\cite{vue-router-guards}, men formularvalidering kræver et tredjepartsbibliotek som VeeValidate. I React er hverken routing eller formularvalidering en del af biblioteket, så begge dele løses med tredjepartsbiblioteker som React Router og React Hook Form. Uanset valg er rollebaseret navigation i frontenden kun en brugeroplevelse, og den egentlige adgangskontrol skal håndhæves i backenden.

\medskip\noindent\textbf{Score:} Angular: 10 \quad React: 7 \quad Vue: 8

\subsection*{State management og API-integration (10\%)}
Frontenden skal holde styr på loginstatus, brugerrolle, opgavelister og saldo. I Angular kan delt state placeres i services via dependency injection og gøres reaktiv med signals~\cite{angular-signals}, mens HttpClient med interceptors kan tilføje JWT-tokens til alle forespørgsler ét sted~\cite{angular-interceptors}. HttpClient bygger dog på RxJS, som tilføjer kompleksitet. Vue har Pinia som officiel løsning til delt state~\cite{pinia}, og React har indbygget \texttt{useState} og \texttt{useContext} suppleret af biblioteker som TanStack Query til server-data. Alle tre løser behovet godt, blot med forskellig grad af indbygget understøttelse.

\medskip\noindent\textbf{Score:} Angular: 8 \quad React: 8 \quad Vue: 8

\subsection*{Type safety og testbarhed (10\%)}
Robusthed er vigtig i flows med betaling, roller og opgavestatus. Angular er bygget til TypeScript, og compileren kan også typetjekke templates~\cite{angular-template-typecheck}. Angular CLI opsætter testmiljøet automatisk med TestBed og Vitest~\cite{angular-testing}. Vue 3 er selv skrevet i TypeScript og har officiel understøttelse af typetjek i templates via \texttt{vue-tsc}~\cite{vue-typescript}, og Vitest anbefales til test~\cite{vue-testing}. React fungerer godt med TypeScript, men typerne vedligeholdes af communityet i pakken \texttt{@types/react}~\cite{react-typescript}, og der findes ingen officiel testopsætning, så gruppen selv skal vælge værktøjer som React Testing Library.

\medskip\noindent\textbf{Score:} Angular: 9 \quad React: 7 \quad Vue: 8

\subsection*{Komponentstruktur og design system (5\%)}
Alle tre er komponentbaserede og understøtter genbrugelige komponenter på tværs af opgaver, rewards, profiler og ratings. Angular har med Angular Material et officielt komponentbibliotek, der sikrer ensartet design og tilgængelighed~\cite{angular-material}. React har det bredeste udvalg af design systems, f.eks. MUI og shadcn/ui, og komponentmodellen er meget fleksibel. Vue har gode biblioteker som Vuetify og PrimeVue, men udvalget er mindre.

\medskip\noindent\textbf{Score:} Angular: 9 \quad React: 9 \quad Vue: 8

\subsection{Scoring}
Scoringen er baseret på en skala fra 1 til 10, hvor 10 er bedst. Den vægtede score er beregnet ud fra kriteriets vægt.

\begin{table}[H]
    \centering
    \renewcommand{\arraystretch}{1.2}
    \begin{tabularx}{\textwidth}{|X|c|c|c|c|}
        \hline
        \rowcolor{gray!30}
        \textbf{Kriterie} & \textbf{Vægt} & \textbf{Angular} & \textbf{React} & \textbf{Vue} \\
        \hline
        Egnethed til projektet & 20\% & 9 & 7 & 7 \\
        \hline
        Maintainability & 10\% & 9 & 7 & 7 \\
        \hline
        Performance & 5\% & 8 & 8 & 8 \\
        \hline
        Kompetencer i gruppen & 10\% & 4 & 5 & 2 \\
        \hline
        Ecosystem/tooling & 5\% & 9 & 10 & 8 \\
        \hline
        Fremtidig skalerbarhed & 5\% & 9 & 9 & 8 \\
        \hline
        Responsivt webdesign og mobilvisning & 10\% & 8 & 8 & 8 \\
        \hline
        Formularer, routing og rollebaseret navigation & 10\% & 10 & 7 & 8 \\
        \hline
        State management og API-integration & 10\% & 8 & 8 & 8 \\
        \hline
        Type safety og testbarhed & 10\% & 9 & 7 & 8 \\
        \hline
        Komponentstruktur og design system & 5\% & 9 & 9 & 8 \\
        \hline
        \rowcolor{gray!30}
        \textbf{Vægtet score} & \textbf{100\%} & \textbf{8,35} & \textbf{7,40} & \textbf{7,10} \\
        \hline
    \end{tabularx}
    \caption{Scoring af frontend-kandidater}
    \label{tab:frontend-scoring}
\end{table}

\subsection*{Beslutning}
På baggrund af analysen vurderes Angular som det mest hensigtsmæssige valg for Echo. Angular scorer højest samlet, fordi frameworket passer godt til en større web-app med mange brugerroller, formularer, routes og API-kald. Samtidig har gruppen markant bedre kompetencer i Angular end i React og Vue, hvilket reducerer risikoen i udviklingsprocessen og gør det mere realistisk at nå en funktionel MVP.\\

Angulars struktur er også en fordel for rapporten, fordi det giver tydelige elementer at dokumentere: komponenter, services, route guards, formularer og dataflow. Valget understøtter derfor både produktets behov og bachelorprojektets behov for en forklarbar teknisk proces.
\end{document}
